\section{RL\_AGENT/EMS\_RL\_agent.py}

\subsection*{Scopo}
Definisce configurazione RL, ambiente Gym offline della microgrid e wrapper agente RL.
Il file centralizza caricamento YAML, costruzione della microgrid, traduzione
dell'azione RL in setpoint batteria e calcolo della reward composita.

\subsection*{Interfacce principali}
\begin{itemize}[leftmargin=1.2cm]
\item \texttt{load\_rl\_agent\_config}: merge config con default robusti.
\item \texttt{OfflineMicrogridRLEnv}: ambiente RL basato su \texttt{gym.Env} con osservazioni/azioni e
  reward shaping.
\item \texttt{EMS\_RL\_Agent}: orchestrazione training/inference con algoritmi SB3
  selezionabili da config (\texttt{PPO} o \texttt{SAC}).
\item Config algoritmo:
  \texttt{rl.algorithm} sceglie l'algoritmo, con blocchi dedicati
  \texttt{rl.ppo} e \texttt{rl.sac}; default retrocompatibile su \texttt{PPO}.
\item Compatibilita PPO con config esterne:
  \texttt{rl.ppo.normalize\_advantage} viene passato al costruttore SB3;
  il campo legacy \texttt{center\_adv} e' supportato come alias.
\item Contract dei builder modello:
  \texttt{\_build\_ppo\_model} e \texttt{\_build\_sac\_model} restituiscono sempre
  l'istanza SB3 costruita; \texttt{build\_model} assegna poi il risultato a
  \texttt{self.model}.
\item Config osservazioni batteria:
  \texttt{soe} e \texttt{soh} sono inclusi nello stato RL; il flag
  \texttt{rl.include\_battery\_power\_limits\_in\_obs} (default \texttt{true}) permette di
  includere/escludere \texttt{max\_production/max\_consumption}.
\item Vincoli opzionali sui flussi energetici:
  \texttt{rl.allow\_grid\_charge\_for\_battery} disabilita la carica da rete quando vale
  \texttt{false}, limitando la carica al solo surplus PV del timestep.
  \texttt{rl.allow\_battery\_export\_to\_grid} puo' limitare la scarica della batteria
  al solo carico locale, evitando export causato dalla batteria.
\item Naming parametri RL coerente con SoE:
  \texttt{randomize\_initial\_SoE}, \texttt{soe\_action\_guard\_band},
  \texttt{coeff\_soe\_violation}, \texttt{soe\_tolerance} (fallback legacy da \texttt{soc\_*} supportato).
\item Reward terms batteria rinominati in log/info da \texttt{soc\_*} a \texttt{soe\_*}
  (es. \texttt{soe\_violation}, \texttt{reward\_term\_soe\_violation}).
\item Modalita reward:
  \texttt{rl.reward.mode: custom} usa la somma pesata dei termini RL;
  \texttt{rl.reward.mode: pymgrid} usa la \texttt{base\_reward}
  restituita da \texttt{microgrid.step}, oppure il delta rispetto a una
  policy di riferimento se la reward relativa e' attiva.
\item Reward relativa opzionale:
  con \texttt{rl.reward.relative\_reference\_enabled: true} il termine economico
  viene sostituito dal delta rispetto a una policy di riferimento
  (\texttt{no\_battery} oppure \texttt{rbc}), con normalizzazione opzionale
  via \texttt{relative\_reference\_normalize}/\texttt{relative\_reference\_eps}.
\item Il termine reward \texttt{soe\_violation} viene letto da \texttt{step\_info} se disponibile,
  con fallback locale su limiti SoE in assenza di campi dedicati.
\item Il clipping dell'azione batteria combina limiti fisici modulo e finestra SoC
  configurata (\texttt{min\_soc/max\_soc}) senza dipendere da un layer BMS.
\item Il termine \texttt{wear\_cost} usa eventuali campi presenti in \texttt{step\_info} o nel
  transition model; se assenti resta nullo.
\item Nel log simulazione, la transition history viene letta se disponibile sul modulo/transition model.
\end{itemize}

\subsection*{Flusso dati}
\begin{itemize}[leftmargin=1.2cm]
\item Carica i dati offline tramite \texttt{RL\_utils\_offline.load\_offline\_timeseries}.
\item In modalita \texttt{scenario.dataset\_mode: pymgrid\_bundle} usa direttamente i file
  \texttt{LoadModule/RenewableModule/GridModule} dello scenario pymgrid, inclusi
  \texttt{price\_buy/price\_sell} reali del \texttt{GridModule}.
\item In modalita CSV usa \texttt{price\_buy/price\_sell} dal dataset se presenti,
  altrimenti fallback alle \texttt{ems.price\_bands}.
\item Istanzia simulatore microgrid.
\item Costruisce osservazioni RL da load, PV, SoE, SoH, prezzi e forecast opzionali.
\item Traduce l'azione RL in richiesta batteria, applica deadband, guard band SoE,
  vincoli opzionali sui flussi energetici e clipping fisico di batteria/rete.
\item Esegue \texttt{microgrid.step}, legge i log modulo e compone reward e diagnostica.
\item La reward finale usata dall'agente dipende da \texttt{rl.reward.mode}:
  in \texttt{custom} e' la somma pesata dei termini RL, in \texttt{pymgrid}
  coincide con la \texttt{base\_reward} del simulatore solo quando
  \texttt{relative\_reference\_enabled=false}.
\item Se la reward relativa e' attiva, il confronto col riferimento viene fatto
  step-by-step e i campi \texttt{reference\_*} vengono esportati in \texttt{info}.
\end{itemize}

\subsection*{Dipendenze critiche}
\begin{itemize}[leftmargin=1.2cm]
\item \texttt{SIMULATOR.microgrid\_simulator.MicrogridSimulator} per costruire la microgrid.
\item \texttt{RL\_utils\_offline.load\_offline\_timeseries} per dataset, prezzi ed emissioni.
\item \texttt{stable\_baselines3}, \texttt{gymnasium}, \texttt{torch}, \texttt{numpy}, \texttt{pandas}, \texttt{yaml}.
\item Battery module base \texttt{pymgrid} per SoE/limiti e transition model opzionale per metriche avanzate.
\end{itemize}

\subsection*{Effetti collaterali}
\begin{itemize}[leftmargin=1.2cm]
\item Legge file YAML di configurazione e dataset CSV/gz dal filesystem.
\item Puo' leggere/scrivere file di observation bounds e statistiche di normalizzazione.
\item Arricchisce \texttt{info} a ogni step con diagnostica dettagliata e popola i log della microgrid.
\item I flag di policy possono ridurre lo spazio di azione effettivamente disponibile all'agente.
\end{itemize}

\subsection*{Rischi e note di manutenzione}
\begin{itemize}[leftmargin=1.2cm]
\item Normalizzazione osservazioni/azioni/reward.
\item Coerenza tra \texttt{rl.algorithm} e checkpoint caricato:
  il loader usa il backend corretto (\texttt{PPO.load} o \texttt{SAC.load}) in base al config.
\item Se cambia la composizione dello stato (es. flag \texttt{include\_battery\_*\_in\_obs}),
  file di bounds/normalizzazione salvati da run precedenti possono non essere compatibili
  (dimensione osservazione diversa).
\item Coerenza tra finestra episodio e lunghezza dataset.
\item Coefficienti reward: possono cambiare fortemente il comportamento dell'agente.
\item I flag sui flussi energetici modificano la policy energetica del problema:
  modelli addestrati senza vincolo possono non restare ottimali con i vincoli attivi, quindi
  e' consigliabile riaddestrare o almeno rivalutare il modello nella nuova configurazione.
\item In scenario-bundle mode, ogni mismatch tra percorsi scenario e file \texttt{.csv.gz}
  produce errore di bootstrap dell'ambiente prima del training.
\end{itemize}
